
\documentclass[11pt]{article}
\usepackage[margin=1in]{geometry}
\usepackage{enumitem}
\usepackage{array}
\usepackage{booktabs}
\setlist{nosep}

\begin{document}

\begin{center}
    {\LARGE Assignment: Pipeline Hazards}\\[0.5em]
    {\large EE161 Homework 6}
\end{center}

\newpage
\section*{Question 1: Dependencies}
Find the data dependencies between all the instructions below. Write your answers in the format of the given example (e.g., ``I2 has a data dependency with I1 in R2'').
\[
\begin{aligned}
\text{I1.}&\ \text{sub\ } R2, R5, R4\\
\text{I2.}&\ \text{add\ } R4, R2, R5\\
\text{I3.}&\ \text{lw\ } R5, 100(R2)\\
\text{I4.}&\ \text{sub\ } R2, R5, R4\\
\text{I5.}&\ \text{sw\ } R2, 101(R2)
\end{aligned}
\]

\newpage
\section*{Question 2: EX forwarding paths}
Add the necessary forwarding paths for the following instructions in the given datapath.

\subsection*{Part A}
\begin{verbatim}
1 add  R3, R3, R4
2 sub  R4, R2, R3
3 add  R4, R5, R5
\end{verbatim}

\subsection*{Part B}
\begin{verbatim}
1 add  R3, R3, R4
2 sub  R4, R2, R4
3 add  R4, R5, R3
\end{verbatim}

\newpage
\section*{Question 3: Performance with stalls}
On a 5-stage MIPS pipeline with forwarding and appropriate NOPs:
\begin{enumerate}[label=\Alph*.]
    \item Fill out the pipeline and draw the forwarding path(s).
    \item How many cycles does it take to execute all three instructions?
\end{enumerate}
\begin{verbatim}
add  R5, R5, R7
lw   R6, 100(R7)
sub  R7, R6, R8
\end{verbatim}

\bigskip
\begin{tabular}{@{}l|ccccccccc@{}}
\textbf{} & \textbf{1} & \textbf{2} & \textbf{3} & \textbf{4} & \textbf{5} & \textbf{6} & \textbf{7} & \textbf{8} & \textbf{9} \\
\midrule
add  R5, R5, R7 & & & & & & & & & \\
lw   R6, 100(R7) & & & & & & & & & \\
sub  R7, R6, R8 & & & & & & & & & \\
\bottomrule
\end{tabular}

\bigskip
Program execution time: \rule{2in}{0.4pt}\, cycles.

\newpage
\section*{Question 4: Load-to-use delay}
The code below was written without regard for delay slots. The pipeline has a double-pumped register file and forwarding.
\begin{verbatim}
lw   R15, 0(R2)
add  R14, R15, R15
lw   R16, 4(R2)
add  R17, R16, R16
\end{verbatim}
\begin{enumerate}[label=\Alph*.]
    \item Identify dependencies.
    \item Insert NOPs to resolve the dependencies.
    \item After adding NOPs, add the instructions and draw the forwarding path(s) in the pipeline.
    \item We would like to keep the pipeline full as much as possible and improve performance. Can we reorder the instructions to avoid NOPs? If so, how?
    \item If the answer to Part D is yes, add the reordered instructions and draw the forwarding path(s) in the pipeline.
\end{enumerate}

\bigskip
\textbf{Load-to-use hint:} the result of \texttt{lw} is only available at the end of MEM, while \texttt{add} needs the value at the start of EXE.

\bigskip
\begin{tabular}{@{}l|ccccccccc@{}}
\textbf{} & \textbf{1} & \textbf{2} & \textbf{3} & \textbf{4} & \textbf{5} & \textbf{6} & \textbf{7} & \textbf{8} & \textbf{9} \\
\midrule
Instruction & & & & & & & & & \\
Instruction & & & & & & & & & \\
Instruction & & & & & & & & & \\
Instruction & & & & & & & & & \\
\bottomrule
\end{tabular}

\newpage
\section*{Question 5: MEM-to-MEM stalls}
Fill in the pipeline for the following code assuming forwarding and double-pumped registers. Show the forwarding path(s) in the pipeline. If there is a dependency that cannot be resolved with forwarding, add bubbles to the pipeline.
\begin{verbatim}
add  R4, R5, R2
lw   R15, 0(R4)
sw   R15, 4(R2)
\end{verbatim}

\bigskip
\begin{tabular}{@{}l|ccccccccc@{}}
\textbf{} & \textbf{1} & \textbf{2} & \textbf{3} & \textbf{4} & \textbf{5} & \textbf{6} & \textbf{7} & \textbf{8} & \textbf{9} \\
\midrule
add  R4, R5, R2  & & & & & & & & & \\
lw   R15, 0(R4)  & & & & & & & & & \\
sw   R15, 4(R2)  & & & & & & & & & \\
\bottomrule
\end{tabular}

\bigskip
Program execution time: \rule{2in}{0.4pt}\, cycles.

\newpage
\section*{Question 6: Branch delay slot}
Parts A and B have two data dependencies (R3 and R4) that can be resolved using forwarding and a double-pumped register file. There is still a control dependency on \texttt{beq}. To solve it, a \texttt{nop} was inserted after \texttt{beq} (see Part B). We want to keep the pipeline full and improve performance.
\begin{enumerate}[label=\Alph*.]
    \item Identify dependencies.
    \item Insert NOPs.
    \item Can we reorder the instructions to avoid NOPs? If so, how?
\end{enumerate}
\begin{verbatim}
add  R3, R3, R4
sub  R4, R2, R3
beq  R0, R2, end
nop
addi R4, R4, 12
end:
\end{verbatim}

\end{document}
